% Khung báo cáo AS3179 (tiếng Việt). Biên dịch: python scripts/04_build_report.py  (pdflatex + T5, cần MiKTeX hoặc TeX Live)
% Các mục đánh số 4.1–4.9 đúng theo đặc tả. Mọi chỗ \TODO{...} do nhóm điền; không có kết luận nào được viết sẵn.
\documentclass[12pt,a4paper]{article}
\usepackage[utf8]{inputenc}
\usepackage[T5]{fontenc}
\usepackage[vietnamese]{babel}
\usepackage{amsmath,amssymb,graphicx,booktabs,longtable,listings,xcolor,hyperref,geometry,float}
\geometry{margin=2.5cm}
\renewcommand{\thesection}{4.\arabic{section}}
\setcounter{secnumdepth}{3}
\newcommand{\TODO}[1]{{\color{red}\textbf{TODO(nhóm):} #1}}
\graphicspath{{../figures/}}
\newcommand{\fig}[2]{\begin{figure}[H]\centering\IfFileExists{../figures/#1.png}{\includegraphics[width=\linewidth]{#1}}{\TODO{thiếu hình #1}}\caption{#2}\end{figure}}
\lstset{basicstyle=\ttfamily\small,breaklines=true,frame=single,language=Python}
\hypersetup{colorlinks=true,linkcolor=black,urlcolor=blue}

\title{Bài tập lớn AS3179 -- Mô hình thống kê\\[4pt]
\large Kiểm chứng Bài 11 (Chuỗi thời gian I: phụ thuộc và mô hình ARMA)\\ và Bài 12 (Chuỗi thời gian II: ARIMA và dự báo) bằng Python}
\author{Khoa Khoa học Ứng dụng, HCMUT\\ \TODO{tên nhóm}}
\date{Hạn nộp bản mềm: 12:00 thứ Hai 30/11/2026}

\begin{document}
\maketitle
\tableofcontents
\newpage

\section{Thông tin chung}
\subsection*{Thành viên nhóm}
\begin{tabular}{@{}llll@{}}\toprule STT & Họ và tên & MSSV & Email \\ \midrule
1 & \TODO{họ tên} & \TODO{MSSV} & \TODO{} \\
2 & \TODO{họ tên} & \TODO{MSSV} & \TODO{} \\
3 & \TODO{họ tên} & \TODO{MSSV} & \TODO{} \\ \bottomrule\end{tabular}

\subsection*{Cặp bài giảng}
Bài giảng thứ nhất: \textbf{Bài 11} -- Chuỗi thời gian I: phụ thuộc và mô hình ARMA. Bài giảng thứ hai: \textbf{Bài 12} -- Chuỗi thời gian II: ARIMA và dự báo.
(Đặc tả còn sót cụm ``hai cặp bài giảng''; nhóm hiểu là một cặp 11 + 12.)

\subsection*{Sáu bộ dữ liệu}
\begin{tabular}{@{}lllll@{}}\toprule Bài & Mã & Ví dụ & Bộ dữ liệu (R \texttt{datasets}) & n \\ \midrule
11 & HU & (i) & \texttt{LakeHuron}: mực nước hồ Huron 1875--1972 & 98 \\
11 & LY & (ii) & \texttt{lynx}: linh miêu bị bẫy 1821--1934, $\log_{10}$ & 114 \\
11 & DR & (iii) & \texttt{EuStockMarkets} (DAX): suất sinh lợi log ngày & 1859 \\
12 & NI & (i) & \texttt{Nile}: lưu lượng sông Nile 1871--1970 & 100 \\
12 & AP & (ii) & \texttt{AirPassengers}: hành khách hàng không 1949--1960, $\log$ & 144 \\
12 & DP & (iii) & \texttt{EuStockMarkets} (DAX): log giá đóng cửa ngày & 1860 \\ \bottomrule\end{tabular}

\section{Tóm tắt nội dung lý thuyết cần kiểm chứng}
Danh sách 28 nội dung (14 của Bài 11, 14 của Bài 12) và một phần mở rộng tùy chọn $\star$. Kiểu: [D] trên dữ liệu thật, [M] mô phỏng có điều khiển (seed 3179).
\begin{longtable}{@{}p{1.1cm}p{8.2cm}p{5cm}@{}}
\toprule
ID & Nội dung & Bộ dữ liệu \\
\midrule
\endhead
11.1 & Chuỗi thời gian, dừng và nhiễu trắng & HU, LY, DR (cả 3) \\
11.2 & Tính chất của hàm tự hiệp phương sai & HU, LY, DR (cả 3) \\
11.3 & Trung bình mẫu khi có phụ thuộc & HU (chính), LY, DR; mô phỏng AR(1) \\
11.4 & ACF mẫu, dải iid, Bartlett và kiểm định Ljung–Box & HU, LY, DR (cả 3) + mô phỏng nhiễu iid \\
11.5 & Quá trình tuyến tính, toán tử lùi, trọng số $\psi$ và ACF lý thuyết & HU, LY (DR: giải thích không áp dụng) \\
11.6 & Mô hình AR(1) & HU (chính); LY, DR: giải thích \\
11.7 & AR(p): tính nhân quả và phương trình Yule–Walker & LY, HU (DR: không áp dụng) \\
11.8 & AR(2): tam giác nhân quả và chu kỳ & LY (chính), HU; DR: không áp dụng \\
11.9 & Mô hình MA(q): ACF cắt cụt, MA(1), khả nghịch & DR, HU (LY: giải thích) + mô phỏng MA(2) \\
11.10 & ARMA(p, q): nhân quả, khả nghịch, nghiệm chung & HU, LY (DR: minh họa thừa tham số) \\
11.11 & Tự tương quan riêng phần (PACF) & HU (đã khử xu thế), LY, DR (cả 3) \\
11.12 & Ước lượng: Yule–Walker, OLS, cực đại hợp lý & LY, HU (+ mô phỏng AR(2)); DR: không áp dụng \\
11.13 & Chọn bậc bằng AIC/BIC và thẩm định cuốn chiếu & LY, HU (DR: lưới ARMA nhỏ) \\
11.14 & Chẩn đoán phần dư và checklist báo cáo ARMA & HU, LY, DR (cả 3) \\
12.1 & Xu thế tất định và xu thế ngẫu nhiên; I(d) & NI, AP, DP (cả 3) \\
12.2 & Bước ngẫu nhiên & DP (mô phỏng, tham số từ DP); NI, AP: giải thích \\
12.3 & Sai phân sửa được gì và làm hỏng gì: sai phân quá mức & NI, DP (AP: xem 12.10) \\
12.4 & Kiểm định nghiệm đơn vị: ADF và KPSS & NI, AP, DP (cả 3) + mô phỏng \\
12.5 & Mô hình ARIMA(p, d, q) và độ trôi & NI, DP (AP: xem 12.10) \\
12.6 & Ảnh hưởng vĩnh viễn của cú sốc: trọng số $\Psi$ & NI, DP (AP: giải thích) \\
12.7 & ARIMA(0,1,1) $\equiv$ làm trơn mũ đơn $\equiv$ mô hình mức cục bộ & NI (chính); DP, AP: giải thích \\
12.8 & Dự báo h bước tối ưu MSE, phương sai sai số và khoảng dự báo & NI, AP, DP (cả 3) + mô phỏng độ phủ \\
12.9 & Dự báo trên thang log: trung vị, trung bình, khoảng chuyển đổi & AP, DP (NI: không biến đổi) \\
12.10 & Sai phân mùa vụ, SARIMA và mô hình airline & AP \\
12.11 & Mốc chuẩn và thước đo độ chính xác & NI, AP, DP (cả 3) \\
12.12 & Thẩm định trên gốc cuốn chiếu và vai trò của bước chọn mô hình & NI, AP, DP (cả 3) \\
12.13 & Hàm mất mát xác định phiếm hàm; kiểm định Diebold–Mariano & AP (chính), NI \\
12.14 & Checklist báo cáo dự báo ARIMA & NI, AP, DP (cả 3) \\
$\star$ & Mở rộng (tùy chọn): hồi quy giả tạo và đồng tích hợp log DAX $\sim$ log SMI & DAX và SMI (phần mở rộng, không thuộc 6 bộ chính) \\
\bottomrule
\end{longtable}

\TODO{thêm cột ``kiểu'' [D]/[M] và số hiệu Định nghĩa/Định lý nếu muốn; không sao chép lại lý thuyết.}

\section{Mô tả và xử lý dữ liệu}
Nguồn, đơn vị và các bước làm sạch: xem \texttt{data/DATA\_SOURCES.md}. Dữ liệu DAX \textbf{không có ngày lịch}: chỉ dùng chỉ số ngày giao dịch $t=1,\dots,1860$.
\subsection*{HU -- \texttt{LakeHuron}}\TODO{nguồn (chép trích dẫn từ \texttt{?LakeHuron}), đơn vị quan sát, biến đáp ứng, tiền xử lý, nhận xét EDA.}\fig{eda_HU_overview}{HU: chuỗi, ACF, PACF.}
\subsection*{LY -- \texttt{lynx}}\TODO{như trên.}\fig{eda_LY_overview}{LY: $\log_{10}$(số lượng), ACF, PACF.}
\subsection*{DR -- DAX, suất sinh lợi}\TODO{như trên.}\fig{eda_DR_overview}{DR: suất sinh lợi log (\%), ACF, PACF.}
\subsection*{NI -- \texttt{Nile}}\TODO{như trên; nêu đổi mức quanh 1899.}\fig{eda_NI_overview}{NI: lưu lượng, ACF, PACF.}
\subsection*{AP -- \texttt{AirPassengers}}\TODO{như trên; log, chia huấn luyện 120 / kiểm định 24.}\fig{eda_AP_overview}{AP: log hành khách, ACF, PACF.}
\subsection*{DP -- DAX, log giá}\TODO{như trên; chia 1500 / 360 ngày giao dịch.}\fig{eda_DP_overview}{DP: log giá, ACF, PACF.}

\section{Kiểm chứng cặp bài giảng thứ nhất (Bài 11)}
Mỗi mục có đủ sáu phần theo đặc tả. Mã nguồn đầy đủ: \texttt{notebooks/bai11\_arma.ipynb}.
% Sinh một lần từ cấu trúc notebook; mỗi mục có đủ 6 tiểu mục theo đặc tả (Bài 11).

\subsection{11.1\quad Chuỗi thời gian, dừng và nhiễu trắng}
\textit{Bộ dữ liệu: HU, LY, DR (cả 3)}

\subsubsection*{Nội dung lý thuyết cần kiểm chứng}
\TODO{viết ngắn gọn, chính xác theo bài giảng, có số hiệu Định nghĩa/Định lý/Mệnh đề (xem ô ``Nội dung lý thuyết'' của notebook).}

\subsubsection*{Cách triển khai}
\TODO{áp dụng trên bộ nào, kiểu [D] hay [M] (xem ô ``Cách triển khai'').}

\subsubsection*{Mã lệnh Python}
\TODO{dán đoạn mã quan trọng dạng văn bản trong khung \texttt{lstlisting} (không chụp ảnh).}

\subsubsection*{Kết quả tính toán}
\TODO{chèn bảng/hình từ \texttt{reports/tables} và \texttt{reports/figures} do notebook sinh ra.}

\subsubsection*{Phân tích và diễn giải}
\TODO{giải thích ý nghĩa thống kê và liên hệ trực tiếp với lý thuyết.}

\subsubsection*{Nhận xét}
\TODO{phù hợp lý thuyết hay không; nếu lệch thì nguyên nhân hợp lý.}

\subsection{11.2\quad Tính chất của hàm tự hiệp phương sai}
\textit{Bộ dữ liệu: HU, LY, DR (cả 3)}

\subsubsection*{Nội dung lý thuyết cần kiểm chứng}
\TODO{viết ngắn gọn, chính xác theo bài giảng, có số hiệu Định nghĩa/Định lý/Mệnh đề (xem ô ``Nội dung lý thuyết'' của notebook).}

\subsubsection*{Cách triển khai}
\TODO{áp dụng trên bộ nào, kiểu [D] hay [M] (xem ô ``Cách triển khai'').}

\subsubsection*{Mã lệnh Python}
\TODO{dán đoạn mã quan trọng dạng văn bản trong khung \texttt{lstlisting} (không chụp ảnh).}

\subsubsection*{Kết quả tính toán}
\TODO{chèn bảng/hình từ \texttt{reports/tables} và \texttt{reports/figures} do notebook sinh ra.}

\subsubsection*{Phân tích và diễn giải}
\TODO{giải thích ý nghĩa thống kê và liên hệ trực tiếp với lý thuyết.}

\subsubsection*{Nhận xét}
\TODO{phù hợp lý thuyết hay không; nếu lệch thì nguyên nhân hợp lý.}

\subsection{11.3\quad Trung bình mẫu khi có phụ thuộc}
\textit{Bộ dữ liệu: HU (chính), LY, DR; mô phỏng AR(1)}

\subsubsection*{Nội dung lý thuyết cần kiểm chứng}
\TODO{viết ngắn gọn, chính xác theo bài giảng, có số hiệu Định nghĩa/Định lý/Mệnh đề (xem ô ``Nội dung lý thuyết'' của notebook).}

\subsubsection*{Cách triển khai}
\TODO{áp dụng trên bộ nào, kiểu [D] hay [M] (xem ô ``Cách triển khai'').}

\subsubsection*{Mã lệnh Python}
\TODO{dán đoạn mã quan trọng dạng văn bản trong khung \texttt{lstlisting} (không chụp ảnh).}

\subsubsection*{Kết quả tính toán}
\TODO{chèn bảng/hình từ \texttt{reports/tables} và \texttt{reports/figures} do notebook sinh ra.}

\subsubsection*{Phân tích và diễn giải}
\TODO{giải thích ý nghĩa thống kê và liên hệ trực tiếp với lý thuyết.}

\subsubsection*{Nhận xét}
\TODO{phù hợp lý thuyết hay không; nếu lệch thì nguyên nhân hợp lý.}

\subsection{11.4\quad ACF mẫu, dải iid, Bartlett và kiểm định Ljung–Box}
\textit{Bộ dữ liệu: HU, LY, DR (cả 3) + mô phỏng nhiễu iid}

\subsubsection*{Nội dung lý thuyết cần kiểm chứng}
\TODO{viết ngắn gọn, chính xác theo bài giảng, có số hiệu Định nghĩa/Định lý/Mệnh đề (xem ô ``Nội dung lý thuyết'' của notebook).}

\subsubsection*{Cách triển khai}
\TODO{áp dụng trên bộ nào, kiểu [D] hay [M] (xem ô ``Cách triển khai'').}

\subsubsection*{Mã lệnh Python}
\TODO{dán đoạn mã quan trọng dạng văn bản trong khung \texttt{lstlisting} (không chụp ảnh).}

\subsubsection*{Kết quả tính toán}
\TODO{chèn bảng/hình từ \texttt{reports/tables} và \texttt{reports/figures} do notebook sinh ra.}

\subsubsection*{Phân tích và diễn giải}
\TODO{giải thích ý nghĩa thống kê và liên hệ trực tiếp với lý thuyết.}

\subsubsection*{Nhận xét}
\TODO{phù hợp lý thuyết hay không; nếu lệch thì nguyên nhân hợp lý.}

\subsection{11.5\quad Quá trình tuyến tính, toán tử lùi, trọng số $\psi$ và ACF lý thuyết}
\textit{Bộ dữ liệu: HU, LY (DR: giải thích không áp dụng)}

\subsubsection*{Nội dung lý thuyết cần kiểm chứng}
\TODO{viết ngắn gọn, chính xác theo bài giảng, có số hiệu Định nghĩa/Định lý/Mệnh đề (xem ô ``Nội dung lý thuyết'' của notebook).}

\subsubsection*{Cách triển khai}
\TODO{áp dụng trên bộ nào, kiểu [D] hay [M] (xem ô ``Cách triển khai'').}

\subsubsection*{Mã lệnh Python}
\TODO{dán đoạn mã quan trọng dạng văn bản trong khung \texttt{lstlisting} (không chụp ảnh).}

\subsubsection*{Kết quả tính toán}
\TODO{chèn bảng/hình từ \texttt{reports/tables} và \texttt{reports/figures} do notebook sinh ra.}

\subsubsection*{Phân tích và diễn giải}
\TODO{giải thích ý nghĩa thống kê và liên hệ trực tiếp với lý thuyết.}

\subsubsection*{Nhận xét}
\TODO{phù hợp lý thuyết hay không; nếu lệch thì nguyên nhân hợp lý.}

\subsection{11.6\quad Mô hình AR(1)}
\textit{Bộ dữ liệu: HU (chính); LY, DR: giải thích}

\subsubsection*{Nội dung lý thuyết cần kiểm chứng}
\TODO{viết ngắn gọn, chính xác theo bài giảng, có số hiệu Định nghĩa/Định lý/Mệnh đề (xem ô ``Nội dung lý thuyết'' của notebook).}

\subsubsection*{Cách triển khai}
\TODO{áp dụng trên bộ nào, kiểu [D] hay [M] (xem ô ``Cách triển khai'').}

\subsubsection*{Mã lệnh Python}
\TODO{dán đoạn mã quan trọng dạng văn bản trong khung \texttt{lstlisting} (không chụp ảnh).}

\subsubsection*{Kết quả tính toán}
\TODO{chèn bảng/hình từ \texttt{reports/tables} và \texttt{reports/figures} do notebook sinh ra.}

\subsubsection*{Phân tích và diễn giải}
\TODO{giải thích ý nghĩa thống kê và liên hệ trực tiếp với lý thuyết.}

\subsubsection*{Nhận xét}
\TODO{phù hợp lý thuyết hay không; nếu lệch thì nguyên nhân hợp lý.}

\subsection{11.7\quad AR(p): tính nhân quả và phương trình Yule–Walker}
\textit{Bộ dữ liệu: LY, HU (DR: không áp dụng)}

\subsubsection*{Nội dung lý thuyết cần kiểm chứng}
\TODO{viết ngắn gọn, chính xác theo bài giảng, có số hiệu Định nghĩa/Định lý/Mệnh đề (xem ô ``Nội dung lý thuyết'' của notebook).}

\subsubsection*{Cách triển khai}
\TODO{áp dụng trên bộ nào, kiểu [D] hay [M] (xem ô ``Cách triển khai'').}

\subsubsection*{Mã lệnh Python}
\TODO{dán đoạn mã quan trọng dạng văn bản trong khung \texttt{lstlisting} (không chụp ảnh).}

\subsubsection*{Kết quả tính toán}
\TODO{chèn bảng/hình từ \texttt{reports/tables} và \texttt{reports/figures} do notebook sinh ra.}

\subsubsection*{Phân tích và diễn giải}
\TODO{giải thích ý nghĩa thống kê và liên hệ trực tiếp với lý thuyết.}

\subsubsection*{Nhận xét}
\TODO{phù hợp lý thuyết hay không; nếu lệch thì nguyên nhân hợp lý.}

\subsection{11.8\quad AR(2): tam giác nhân quả và chu kỳ}
\textit{Bộ dữ liệu: LY (chính), HU; DR: không áp dụng}

\subsubsection*{Nội dung lý thuyết cần kiểm chứng}
\TODO{viết ngắn gọn, chính xác theo bài giảng, có số hiệu Định nghĩa/Định lý/Mệnh đề (xem ô ``Nội dung lý thuyết'' của notebook).}

\subsubsection*{Cách triển khai}
\TODO{áp dụng trên bộ nào, kiểu [D] hay [M] (xem ô ``Cách triển khai'').}

\subsubsection*{Mã lệnh Python}
\TODO{dán đoạn mã quan trọng dạng văn bản trong khung \texttt{lstlisting} (không chụp ảnh).}

\subsubsection*{Kết quả tính toán}
\TODO{chèn bảng/hình từ \texttt{reports/tables} và \texttt{reports/figures} do notebook sinh ra.}

\subsubsection*{Phân tích và diễn giải}
\TODO{giải thích ý nghĩa thống kê và liên hệ trực tiếp với lý thuyết.}

\subsubsection*{Nhận xét}
\TODO{phù hợp lý thuyết hay không; nếu lệch thì nguyên nhân hợp lý.}

\subsection{11.9\quad Mô hình MA(q): ACF cắt cụt, MA(1), khả nghịch}
\textit{Bộ dữ liệu: DR, HU (LY: giải thích) + mô phỏng MA(2)}

\subsubsection*{Nội dung lý thuyết cần kiểm chứng}
\TODO{viết ngắn gọn, chính xác theo bài giảng, có số hiệu Định nghĩa/Định lý/Mệnh đề (xem ô ``Nội dung lý thuyết'' của notebook).}

\subsubsection*{Cách triển khai}
\TODO{áp dụng trên bộ nào, kiểu [D] hay [M] (xem ô ``Cách triển khai'').}

\subsubsection*{Mã lệnh Python}
\TODO{dán đoạn mã quan trọng dạng văn bản trong khung \texttt{lstlisting} (không chụp ảnh).}

\subsubsection*{Kết quả tính toán}
\TODO{chèn bảng/hình từ \texttt{reports/tables} và \texttt{reports/figures} do notebook sinh ra.}

\subsubsection*{Phân tích và diễn giải}
\TODO{giải thích ý nghĩa thống kê và liên hệ trực tiếp với lý thuyết.}

\subsubsection*{Nhận xét}
\TODO{phù hợp lý thuyết hay không; nếu lệch thì nguyên nhân hợp lý.}

\subsection{11.10\quad ARMA(p, q): nhân quả, khả nghịch, nghiệm chung}
\textit{Bộ dữ liệu: HU, LY (DR: minh họa thừa tham số)}

\subsubsection*{Nội dung lý thuyết cần kiểm chứng}
\TODO{viết ngắn gọn, chính xác theo bài giảng, có số hiệu Định nghĩa/Định lý/Mệnh đề (xem ô ``Nội dung lý thuyết'' của notebook).}

\subsubsection*{Cách triển khai}
\TODO{áp dụng trên bộ nào, kiểu [D] hay [M] (xem ô ``Cách triển khai'').}

\subsubsection*{Mã lệnh Python}
\TODO{dán đoạn mã quan trọng dạng văn bản trong khung \texttt{lstlisting} (không chụp ảnh).}

\subsubsection*{Kết quả tính toán}
\TODO{chèn bảng/hình từ \texttt{reports/tables} và \texttt{reports/figures} do notebook sinh ra.}

\subsubsection*{Phân tích và diễn giải}
\TODO{giải thích ý nghĩa thống kê và liên hệ trực tiếp với lý thuyết.}

\subsubsection*{Nhận xét}
\TODO{phù hợp lý thuyết hay không; nếu lệch thì nguyên nhân hợp lý.}

\subsection{11.11\quad Tự tương quan riêng phần (PACF)}
\textit{Bộ dữ liệu: HU (đã khử xu thế), LY, DR (cả 3)}

\subsubsection*{Nội dung lý thuyết cần kiểm chứng}
\TODO{viết ngắn gọn, chính xác theo bài giảng, có số hiệu Định nghĩa/Định lý/Mệnh đề (xem ô ``Nội dung lý thuyết'' của notebook).}

\subsubsection*{Cách triển khai}
\TODO{áp dụng trên bộ nào, kiểu [D] hay [M] (xem ô ``Cách triển khai'').}

\subsubsection*{Mã lệnh Python}
\TODO{dán đoạn mã quan trọng dạng văn bản trong khung \texttt{lstlisting} (không chụp ảnh).}

\subsubsection*{Kết quả tính toán}
\TODO{chèn bảng/hình từ \texttt{reports/tables} và \texttt{reports/figures} do notebook sinh ra.}

\subsubsection*{Phân tích và diễn giải}
\TODO{giải thích ý nghĩa thống kê và liên hệ trực tiếp với lý thuyết.}

\subsubsection*{Nhận xét}
\TODO{phù hợp lý thuyết hay không; nếu lệch thì nguyên nhân hợp lý.}

\subsection{11.12\quad Ước lượng: Yule–Walker, OLS, cực đại hợp lý}
\textit{Bộ dữ liệu: LY, HU (+ mô phỏng AR(2)); DR: không áp dụng}

\subsubsection*{Nội dung lý thuyết cần kiểm chứng}
\TODO{viết ngắn gọn, chính xác theo bài giảng, có số hiệu Định nghĩa/Định lý/Mệnh đề (xem ô ``Nội dung lý thuyết'' của notebook).}

\subsubsection*{Cách triển khai}
\TODO{áp dụng trên bộ nào, kiểu [D] hay [M] (xem ô ``Cách triển khai'').}

\subsubsection*{Mã lệnh Python}
\TODO{dán đoạn mã quan trọng dạng văn bản trong khung \texttt{lstlisting} (không chụp ảnh).}

\subsubsection*{Kết quả tính toán}
\TODO{chèn bảng/hình từ \texttt{reports/tables} và \texttt{reports/figures} do notebook sinh ra.}

\subsubsection*{Phân tích và diễn giải}
\TODO{giải thích ý nghĩa thống kê và liên hệ trực tiếp với lý thuyết.}

\subsubsection*{Nhận xét}
\TODO{phù hợp lý thuyết hay không; nếu lệch thì nguyên nhân hợp lý.}

\subsection{11.13\quad Chọn bậc bằng AIC/BIC và thẩm định cuốn chiếu}
\textit{Bộ dữ liệu: LY, HU (DR: lưới ARMA nhỏ)}

\subsubsection*{Nội dung lý thuyết cần kiểm chứng}
\TODO{viết ngắn gọn, chính xác theo bài giảng, có số hiệu Định nghĩa/Định lý/Mệnh đề (xem ô ``Nội dung lý thuyết'' của notebook).}

\subsubsection*{Cách triển khai}
\TODO{áp dụng trên bộ nào, kiểu [D] hay [M] (xem ô ``Cách triển khai'').}

\subsubsection*{Mã lệnh Python}
\TODO{dán đoạn mã quan trọng dạng văn bản trong khung \texttt{lstlisting} (không chụp ảnh).}

\subsubsection*{Kết quả tính toán}
\TODO{chèn bảng/hình từ \texttt{reports/tables} và \texttt{reports/figures} do notebook sinh ra.}

\subsubsection*{Phân tích và diễn giải}
\TODO{giải thích ý nghĩa thống kê và liên hệ trực tiếp với lý thuyết.}

\subsubsection*{Nhận xét}
\TODO{phù hợp lý thuyết hay không; nếu lệch thì nguyên nhân hợp lý.}

\subsection{11.14\quad Chẩn đoán phần dư và checklist báo cáo ARMA}
\textit{Bộ dữ liệu: HU, LY, DR (cả 3)}

\subsubsection*{Nội dung lý thuyết cần kiểm chứng}
\TODO{viết ngắn gọn, chính xác theo bài giảng, có số hiệu Định nghĩa/Định lý/Mệnh đề (xem ô ``Nội dung lý thuyết'' của notebook).}

\subsubsection*{Cách triển khai}
\TODO{áp dụng trên bộ nào, kiểu [D] hay [M] (xem ô ``Cách triển khai'').}

\subsubsection*{Mã lệnh Python}
\TODO{dán đoạn mã quan trọng dạng văn bản trong khung \texttt{lstlisting} (không chụp ảnh).}

\subsubsection*{Kết quả tính toán}
\TODO{chèn bảng/hình từ \texttt{reports/tables} và \texttt{reports/figures} do notebook sinh ra.}

\subsubsection*{Phân tích và diễn giải}
\TODO{giải thích ý nghĩa thống kê và liên hệ trực tiếp với lý thuyết.}

\subsubsection*{Nhận xét}
\TODO{phù hợp lý thuyết hay không; nếu lệch thì nguyên nhân hợp lý.}


\section{Kiểm chứng cặp bài giảng thứ hai (Bài 12)}
Mã nguồn đầy đủ: \texttt{notebooks/bai12\_arima\_du\_bao.ipynb}.
% Sinh từ cấu trúc notebook; mỗi mục có đủ 6 tiểu mục theo đặc tả (Bài 12). Hình/bảng chèn tự động theo tên fig_b*/tab_b*.

\subsection{12.1\quad Xu thế tất định và xu thế ngẫu nhiên; I(d)}
\textit{Bộ dữ liệu: NI, AP, DP (cả 3)}

\subsubsection*{Nội dung lý thuyết cần kiểm chứng}
\TODO{viết ngắn gọn, chính xác theo bài giảng, có số hiệu Định nghĩa/Định lý/Mệnh đề (xem ô ``Nội dung lý thuyết'' của notebook).}

\subsubsection*{Cách triển khai}
\TODO{áp dụng trên bộ nào, kiểu [D] hay [M] (xem ô ``Cách triển khai'').}

\subsubsection*{Mã lệnh Python}
\TODO{dán đoạn mã quan trọng dạng văn bản trong khung \texttt{lstlisting} (không chụp ảnh).}

\subsubsection*{Kết quả tính toán}
\TODO{chọn những bảng/hình dưới đây cần giữ trong báo cáo (có thể bỏ bớt), sắp xếp và viết câu dẫn.}

\input{\tabdir/tab_b12_01_levels_vs_differences.tex}

\figx{fig_b12_01_levels_vs_differences}{Chuỗi, ACF của chuỗi gốc và ACF của sai phân bậc một cho NI, AP (log, huấn luyện), DP (log)}

\subsubsection*{Phân tích và diễn giải}
\TODO{giải thích ý nghĩa thống kê và liên hệ trực tiếp với lý thuyết.}

\subsubsection*{Nhận xét}
\TODO{phù hợp lý thuyết hay không; nếu lệch thì nguyên nhân hợp lý.}

\subsection{12.2\quad Bước ngẫu nhiên}
\textit{Bộ dữ liệu: DP (mô phỏng, tham số từ DP); NI, AP: giải thích}

\subsubsection*{Nội dung lý thuyết cần kiểm chứng}
\TODO{viết ngắn gọn, chính xác theo bài giảng, có số hiệu Định nghĩa/Định lý/Mệnh đề (xem ô ``Nội dung lý thuyết'' của notebook).}

\subsubsection*{Cách triển khai}
\TODO{áp dụng trên bộ nào, kiểu [D] hay [M] (xem ô ``Cách triển khai'').}

\subsubsection*{Mã lệnh Python}
\TODO{dán đoạn mã quan trọng dạng văn bản trong khung \texttt{lstlisting} (không chụp ảnh).}

\subsubsection*{Kết quả tính toán}
\TODO{chọn những bảng/hình dưới đây cần giữ trong báo cáo (có thể bỏ bớt), sắp xếp và viết câu dẫn.}

\input{\tabdir/tab_b12_02_rw_correlation.tex}

\input{\tabdir/tab_b12_02_rw_sample_acf.tex}

\input{\tabdir/tab_b12_02_rw_variance.tex}

\input{\tabdir/tab_b12_02_spurious_regression.tex}

\figx{fig_b12_02_rw_variance}{Bước ngẫu nhiên mô phỏng: $\operatorname{Var}Y_t$ so với $t\hat\sigma^2$ (trái), 50 đường mẫu (phải)}

\subsubsection*{Phân tích và diễn giải}
\TODO{giải thích ý nghĩa thống kê và liên hệ trực tiếp với lý thuyết.}

\subsubsection*{Nhận xét}
\TODO{phù hợp lý thuyết hay không; nếu lệch thì nguyên nhân hợp lý.}

\subsection{12.3\quad Sai phân sửa được gì và làm hỏng gì: sai phân quá mức}
\textit{Bộ dữ liệu: NI, DP (AP: xem 12.10)}

\subsubsection*{Nội dung lý thuyết cần kiểm chứng}
\TODO{viết ngắn gọn, chính xác theo bài giảng, có số hiệu Định nghĩa/Định lý/Mệnh đề (xem ô ``Nội dung lý thuyết'' của notebook).}

\subsubsection*{Cách triển khai}
\TODO{áp dụng trên bộ nào, kiểu [D] hay [M] (xem ô ``Cách triển khai'').}

\subsubsection*{Mã lệnh Python}
\TODO{dán đoạn mã quan trọng dạng văn bản trong khung \texttt{lstlisting} (không chụp ảnh).}

\subsubsection*{Kết quả tính toán}
\TODO{chọn những bảng/hình dưới đây cần giữ trong báo cáo (có thể bỏ bớt), sắp xếp và viết câu dẫn.}

\input{\tabdir/tab_b12_03_cost_simulation.tex}

\input{\tabdir/tab_b12_03_overdifferencing.tex}

\subsubsection*{Phân tích và diễn giải}
\TODO{giải thích ý nghĩa thống kê và liên hệ trực tiếp với lý thuyết.}

\subsubsection*{Nhận xét}
\TODO{phù hợp lý thuyết hay không; nếu lệch thì nguyên nhân hợp lý.}

\subsection{12.4\quad Kiểm định nghiệm đơn vị: ADF và KPSS}
\textit{Bộ dữ liệu: NI, AP, DP (cả 3) + mô phỏng}

\subsubsection*{Nội dung lý thuyết cần kiểm chứng}
\TODO{viết ngắn gọn, chính xác theo bài giảng, có số hiệu Định nghĩa/Định lý/Mệnh đề (xem ô ``Nội dung lý thuyết'' của notebook).}

\subsubsection*{Cách triển khai}
\TODO{áp dụng trên bộ nào, kiểu [D] hay [M] (xem ô ``Cách triển khai'').}

\subsubsection*{Mã lệnh Python}
\TODO{dán đoạn mã quan trọng dạng văn bản trong khung \texttt{lstlisting} (không chụp ảnh).}

\subsubsection*{Kết quả tính toán}
\TODO{chọn những bảng/hình dưới đây cần giữ trong báo cáo (có thể bỏ bớt), sắp xếp và viết câu dẫn.}

\input{\tabdir/tab_b12_04_adf_power.tex}

\input{\tabdir/tab_b12_04_null_distribution.tex}

\input{\tabdir/tab_b12_04_unit_root_tests.tex}

\figx{fig_b12_04_dickey_fuller_null}{Phân phối mô phỏng của $\tau$ dưới H$_0$ so với $N(0,1)$ và giá trị tới hạn 5\%}

\subsubsection*{Phân tích và diễn giải}
\TODO{giải thích ý nghĩa thống kê và liên hệ trực tiếp với lý thuyết.}

\subsubsection*{Nhận xét}
\TODO{phù hợp lý thuyết hay không; nếu lệch thì nguyên nhân hợp lý.}

\subsection{12.5\quad Mô hình ARIMA(p, d, q) và độ trôi}
\textit{Bộ dữ liệu: NI, DP (AP: xem 12.10)}

\subsubsection*{Nội dung lý thuyết cần kiểm chứng}
\TODO{viết ngắn gọn, chính xác theo bài giảng, có số hiệu Định nghĩa/Định lý/Mệnh đề (xem ô ``Nội dung lý thuyết'' của notebook).}

\subsubsection*{Cách triển khai}
\TODO{áp dụng trên bộ nào, kiểu [D] hay [M] (xem ô ``Cách triển khai'').}

\subsubsection*{Mã lệnh Python}
\TODO{dán đoạn mã quan trọng dạng văn bản trong khung \texttt{lstlisting} (không chụp ảnh).}

\subsubsection*{Kết quả tính toán}
\TODO{chọn những bảng/hình dưới đây cần giữ trong báo cáo (có thể bỏ bớt), sắp xếp và viết câu dẫn.}

\input{\tabdir/tab_b12_05_arima_params.tex}

\input{\tabdir/tab_b12_05_fit_stats.tex}

\input{\tabdir/tab_b12_05_rw_manual_vs_arima.tex}

\figx{fig_b12_05_drift_forecast}{DP: dự báo 100 ngày giao dịch của log DAX có và không có độ trôi}

\subsubsection*{Phân tích và diễn giải}
\TODO{giải thích ý nghĩa thống kê và liên hệ trực tiếp với lý thuyết.}

\subsubsection*{Nhận xét}
\TODO{phù hợp lý thuyết hay không; nếu lệch thì nguyên nhân hợp lý.}

\subsection{12.6\quad Ảnh hưởng vĩnh viễn của cú sốc: trọng số $\Psi$}
\textit{Bộ dữ liệu: NI, DP (AP: giải thích)}

\subsubsection*{Nội dung lý thuyết cần kiểm chứng}
\TODO{viết ngắn gọn, chính xác theo bài giảng, có số hiệu Định nghĩa/Định lý/Mệnh đề (xem ô ``Nội dung lý thuyết'' của notebook).}

\subsubsection*{Cách triển khai}
\TODO{áp dụng trên bộ nào, kiểu [D] hay [M] (xem ô ``Cách triển khai'').}

\subsubsection*{Mã lệnh Python}
\TODO{dán đoạn mã quan trọng dạng văn bản trong khung \texttt{lstlisting} (không chụp ảnh).}

\subsubsection*{Kết quả tính toán}
\TODO{chọn những bảng/hình dưới đây cần giữ trong báo cáo (có thể bỏ bớt), sắp xếp và viết câu dẫn.}

\input{\tabdir/tab_b12_06_psi_weights.tex}

\figx{fig_b12_06_psi_weights}{Ảnh hưởng của một cú sốc đơn vị lên $Y_{t+k}$ ($\Psi_k$)}

\subsubsection*{Phân tích và diễn giải}
\TODO{giải thích ý nghĩa thống kê và liên hệ trực tiếp với lý thuyết.}

\subsubsection*{Nhận xét}
\TODO{phù hợp lý thuyết hay không; nếu lệch thì nguyên nhân hợp lý.}

\subsection{12.7\quad ARIMA(0,1,1) $\equiv$ làm trơn mũ đơn $\equiv$ mô hình mức cục bộ}
\textit{Bộ dữ liệu: NI (chính); DP, AP: giải thích}

\subsubsection*{Nội dung lý thuyết cần kiểm chứng}
\TODO{viết ngắn gọn, chính xác theo bài giảng, có số hiệu Định nghĩa/Định lý/Mệnh đề (xem ô ``Nội dung lý thuyết'' của notebook).}

\subsubsection*{Cách triển khai}
\TODO{áp dụng trên bộ nào, kiểu [D] hay [M] (xem ô ``Cách triển khai'').}

\subsubsection*{Mã lệnh Python}
\TODO{dán đoạn mã quan trọng dạng văn bản trong khung \texttt{lstlisting} (không chụp ảnh).}

\subsubsection*{Kết quả tính toán}
\TODO{chọn những bảng/hình dưới đây cần giữ trong báo cáo (có thể bỏ bớt), sắp xếp và viết câu dẫn.}

\input{\tabdir/tab_b12_07_local_level.tex}

\input{\tabdir/tab_b12_07_ses_forecasts.tex}

\subsubsection*{Phân tích và diễn giải}
\TODO{giải thích ý nghĩa thống kê và liên hệ trực tiếp với lý thuyết.}

\subsubsection*{Nhận xét}
\TODO{phù hợp lý thuyết hay không; nếu lệch thì nguyên nhân hợp lý.}

\subsection{12.8\quad Dự báo h bước tối ưu MSE, phương sai sai số và khoảng dự báo}
\textit{Bộ dữ liệu: NI, AP, DP (cả 3) + mô phỏng độ phủ}

\subsubsection*{Nội dung lý thuyết cần kiểm chứng}
\TODO{viết ngắn gọn, chính xác theo bài giảng, có số hiệu Định nghĩa/Định lý/Mệnh đề (xem ô ``Nội dung lý thuyết'' của notebook).}

\subsubsection*{Cách triển khai}
\TODO{áp dụng trên bộ nào, kiểu [D] hay [M] (xem ô ``Cách triển khai'').}

\subsubsection*{Mã lệnh Python}
\TODO{dán đoạn mã quan trọng dạng văn bản trong khung \texttt{lstlisting} (không chụp ảnh).}

\subsubsection*{Kết quả tính toán}
\TODO{chọn những bảng/hình dưới đây cần giữ trong báo cáo (có thể bỏ bớt), sắp xếp và viết câu dẫn.}

\input{\tabdir/tab_b12_08_coverage.tex}

\input{\tabdir/tab_b12_08_forecast_variance.tex}

\figx{fig_b12_08_ni_forecast}{NI: dự báo ARIMA(0,1,1) và khoảng dự báo 95\% cho 20 năm tiếp theo}

\figx{fig_b12_08_se_growth}{Sai số chuẩn của dự báo $h$ bước ($\sigma=1$): $d=0$ so với $d=1$}

\subsubsection*{Phân tích và diễn giải}
\TODO{giải thích ý nghĩa thống kê và liên hệ trực tiếp với lý thuyết.}

\subsubsection*{Nhận xét}
\TODO{phù hợp lý thuyết hay không; nếu lệch thì nguyên nhân hợp lý.}

\subsection{12.9\quad Dự báo trên thang log: trung vị, trung bình, khoảng chuyển đổi}
\textit{Bộ dữ liệu: AP, DP (NI: không biến đổi)}

\subsubsection*{Nội dung lý thuyết cần kiểm chứng}
\TODO{viết ngắn gọn, chính xác theo bài giảng, có số hiệu Định nghĩa/Định lý/Mệnh đề (xem ô ``Nội dung lý thuyết'' của notebook).}

\subsubsection*{Cách triển khai}
\TODO{áp dụng trên bộ nào, kiểu [D] hay [M] (xem ô ``Cách triển khai'').}

\subsubsection*{Mã lệnh Python}
\TODO{dán đoạn mã quan trọng dạng văn bản trong khung \texttt{lstlisting} (không chụp ảnh).}

\subsubsection*{Kết quả tính toán}
\TODO{chọn những bảng/hình dưới đây cần giữ trong báo cáo (có thể bỏ bớt), sắp xếp và viết câu dẫn.}

\input{\tabdir/tab_b12_09_ap_backtransform.tex}

\input{\tabdir/tab_b12_09_dp_forecast.tex}

\figx{fig_b12_09_ap_forecast}{AP: dự báo airline 1959--1960 trên thang gốc (trung vị, khoảng 95\%) và giá trị thực}

\subsubsection*{Phân tích và diễn giải}
\TODO{giải thích ý nghĩa thống kê và liên hệ trực tiếp với lý thuyết.}

\subsubsection*{Nhận xét}
\TODO{phù hợp lý thuyết hay không; nếu lệch thì nguyên nhân hợp lý.}

\subsection{12.10\quad Sai phân mùa vụ, SARIMA và mô hình airline}
\textit{Bộ dữ liệu: AP}

\subsubsection*{Nội dung lý thuyết cần kiểm chứng}
\TODO{viết ngắn gọn, chính xác theo bài giảng, có số hiệu Định nghĩa/Định lý/Mệnh đề (xem ô ``Nội dung lý thuyết'' của notebook).}

\subsubsection*{Cách triển khai}
\TODO{áp dụng trên bộ nào, kiểu [D] hay [M] (xem ô ``Cách triển khai'').}

\subsubsection*{Mã lệnh Python}
\TODO{dán đoạn mã quan trọng dạng văn bản trong khung \texttt{lstlisting} (không chụp ảnh).}

\subsubsection*{Kết quả tính toán}
\TODO{chọn những bảng/hình dưới đây cần giữ trong báo cáo (có thể bỏ bớt), sắp xếp và viết câu dẫn.}

\input{\tabdir/tab_b12_10_airline_acf.tex}

\input{\tabdir/tab_b12_10_sd.tex}

\input{\tabdir/tab_b12_10_seasonality_models.tex}

\figx{fig_b12_10_airline_acf}{ACF của $\nabla\nabla_{12}$ log AP (que) và ACF lý thuyết của mô hình airline (chấm)}

\subsubsection*{Phân tích và diễn giải}
\TODO{giải thích ý nghĩa thống kê và liên hệ trực tiếp với lý thuyết.}

\subsubsection*{Nhận xét}
\TODO{phù hợp lý thuyết hay không; nếu lệch thì nguyên nhân hợp lý.}

\subsection{12.11\quad Mốc chuẩn và thước đo độ chính xác}
\textit{Bộ dữ liệu: NI, AP, DP (cả 3)}

\subsubsection*{Nội dung lý thuyết cần kiểm chứng}
\TODO{viết ngắn gọn, chính xác theo bài giảng, có số hiệu Định nghĩa/Định lý/Mệnh đề (xem ô ``Nội dung lý thuyết'' của notebook).}

\subsubsection*{Cách triển khai}
\TODO{áp dụng trên bộ nào, kiểu [D] hay [M] (xem ô ``Cách triển khai'').}

\subsubsection*{Mã lệnh Python}
\TODO{dán đoạn mã quan trọng dạng văn bản trong khung \texttt{lstlisting} (không chụp ảnh).}

\subsubsection*{Kết quả tính toán}
\TODO{chọn những bảng/hình dưới đây cần giữ trong báo cáo (có thể bỏ bớt), sắp xếp và viết câu dẫn.}

\input{\tabdir/tab_b12_11_ap_benchmarks.tex}

\input{\tabdir/tab_b12_11_dp_multistep.tex}

\input{\tabdir/tab_b12_11_dp_one_step.tex}

\input{\tabdir/tab_b12_11_ni_benchmarks.tex}

\subsubsection*{Phân tích và diễn giải}
\TODO{giải thích ý nghĩa thống kê và liên hệ trực tiếp với lý thuyết.}

\subsubsection*{Nhận xét}
\TODO{phù hợp lý thuyết hay không; nếu lệch thì nguyên nhân hợp lý.}

\subsection{12.12\quad Thẩm định trên gốc cuốn chiếu và vai trò của bước chọn mô hình}
\textit{Bộ dữ liệu: NI, AP, DP (cả 3)}

\subsubsection*{Nội dung lý thuyết cần kiểm chứng}
\TODO{viết ngắn gọn, chính xác theo bài giảng, có số hiệu Định nghĩa/Định lý/Mệnh đề (xem ô ``Nội dung lý thuyết'' của notebook).}

\subsubsection*{Cách triển khai}
\TODO{áp dụng trên bộ nào, kiểu [D] hay [M] (xem ô ``Cách triển khai'').}

\subsubsection*{Mã lệnh Python}
\TODO{dán đoạn mã quan trọng dạng văn bản trong khung \texttt{lstlisting} (không chụp ảnh).}

\subsubsection*{Kết quả tính toán}
\TODO{chọn những bảng/hình dưới đây cần giữ trong báo cáo (có thể bỏ bớt), sắp xếp và viết câu dẫn.}

\input{\tabdir/tab_b12_12_ap_one_step.tex}

\input{\tabdir/tab_b12_12_ap_rolling_by_horizon.tex}

\input{\tabdir/tab_b12_12_dp_rolling.tex}

\input{\tabdir/tab_b12_12_ni_rolling.tex}

\figx{fig_b12_12_ap_rmse_by_horizon}{AP: RMSE gốc cuốn chiếu theo chân trời, airline so với naive mùa vụ}

\subsubsection*{Phân tích và diễn giải}
\TODO{giải thích ý nghĩa thống kê và liên hệ trực tiếp với lý thuyết.}

\subsubsection*{Nhận xét}
\TODO{phù hợp lý thuyết hay không; nếu lệch thì nguyên nhân hợp lý.}

\subsection{12.13\quad Hàm mất mát xác định phiếm hàm; kiểm định Diebold–Mariano}
\textit{Bộ dữ liệu: AP (chính), NI}

\subsubsection*{Nội dung lý thuyết cần kiểm chứng}
\TODO{viết ngắn gọn, chính xác theo bài giảng, có số hiệu Định nghĩa/Định lý/Mệnh đề (xem ô ``Nội dung lý thuyết'' của notebook).}

\subsubsection*{Cách triển khai}
\TODO{áp dụng trên bộ nào, kiểu [D] hay [M] (xem ô ``Cách triển khai'').}

\subsubsection*{Mã lệnh Python}
\TODO{dán đoạn mã quan trọng dạng văn bản trong khung \texttt{lstlisting} (không chụp ảnh).}

\subsubsection*{Kết quả tính toán}
\TODO{chọn những bảng/hình dưới đây cần giữ trong báo cáo (có thể bỏ bớt), sắp xếp và viết câu dẫn.}

\input{\tabdir/tab_b12_13_diebold_mariano.tex}

\input{\tabdir/tab_b12_13_loss_functional.tex}

\input{\tabdir/tab_b12_13_ranking.tex}

\subsubsection*{Phân tích và diễn giải}
\TODO{giải thích ý nghĩa thống kê và liên hệ trực tiếp với lý thuyết.}

\subsubsection*{Nhận xét}
\TODO{phù hợp lý thuyết hay không; nếu lệch thì nguyên nhân hợp lý.}

\subsection{12.14\quad Checklist báo cáo dự báo ARIMA}
\textit{Bộ dữ liệu: NI, AP, DP (cả 3)}

\subsubsection*{Nội dung lý thuyết cần kiểm chứng}
\TODO{viết ngắn gọn, chính xác theo bài giảng, có số hiệu Định nghĩa/Định lý/Mệnh đề (xem ô ``Nội dung lý thuyết'' của notebook).}

\subsubsection*{Cách triển khai}
\TODO{áp dụng trên bộ nào, kiểu [D] hay [M] (xem ô ``Cách triển khai'').}

\subsubsection*{Mã lệnh Python}
\TODO{dán đoạn mã quan trọng dạng văn bản trong khung \texttt{lstlisting} (không chụp ảnh).}

\subsubsection*{Kết quả tính toán}
\TODO{chọn những bảng/hình dưới đây cần giữ trong báo cáo (có thể bỏ bớt), sắp xếp và viết câu dẫn.}

\input{\tabdir/tab_b12_14_checklist.tex}

\figx{fig_b12_14_residual_acf}{ACF của phần dư (trái) và của phần dư bình phương (phải) cho NI, AP, DP}

\subsubsection*{Phân tích và diễn giải}
\TODO{giải thích ý nghĩa thống kê và liên hệ trực tiếp với lý thuyết.}

\subsubsection*{Nhận xét}
\TODO{phù hợp lý thuyết hay không; nếu lệch thì nguyên nhân hợp lý.}

\subsection{$\star$\quad Mở rộng (tùy chọn): hồi quy giả tạo và đồng tích hợp log DAX $\sim$ log SMI}
\textit{Bộ dữ liệu: DAX và SMI (phần mở rộng, không thuộc 6 bộ chính)}

\subsubsection*{Nội dung lý thuyết cần kiểm chứng}
\TODO{viết ngắn gọn, chính xác theo bài giảng, có số hiệu Định nghĩa/Định lý/Mệnh đề (xem ô ``Nội dung lý thuyết'' của notebook).}

\subsubsection*{Cách triển khai}
\TODO{áp dụng trên bộ nào, kiểu [D] hay [M] (xem ô ``Cách triển khai'').}

\subsubsection*{Mã lệnh Python}
\TODO{dán đoạn mã quan trọng dạng văn bản trong khung \texttt{lstlisting} (không chụp ảnh).}

\subsubsection*{Kết quả tính toán}
\TODO{chọn những bảng/hình dưới đây cần giữ trong báo cáo (có thể bỏ bớt), sắp xếp và viết câu dẫn.}

\input{\tabdir/tab_b12_ext_cointegration.tex}

\subsubsection*{Phân tích và diễn giải}
\TODO{giải thích ý nghĩa thống kê và liên hệ trực tiếp với lý thuyết.}

\subsubsection*{Nhận xét}
\TODO{phù hợp lý thuyết hay không; nếu lệch thì nguyên nhân hợp lý.}


\section{So sánh và thảo luận}
\subsection*{Mức độ phù hợp của kết quả thực nghiệm với lý thuyết}\TODO{tổng hợp từ bảng cuối hai notebook.}
\subsection*{Sự khác biệt giữa các bộ dữ liệu}\TODO{}
\subsection*{Ảnh hưởng của đặc điểm dữ liệu đến kết quả mô hình}\TODO{xu thế, đổi mức, mùa vụ, cụm biến động\dots}
\subsection*{Những giả thiết được đáp ứng hoặc không được đáp ứng}\TODO{}
\subsection*{Hạn chế và vấn đề phát sinh}\TODO{gồm các chỗ lệch giữa R và Python đã ghi ở \texttt{reports/tables/lecture\_reference\_check.md}.}

\section{Kết luận}
\TODO{tóm tắt kết quả quan trọng từ hai bài giảng và sáu bộ dữ liệu, nhận xét chính.}

\section{Đóng góp của các thành viên}
\begin{tabular}{@{}p{4cm}p{9cm}p{2cm}@{}}\toprule Thành viên & Công việc cụ thể & Tỷ lệ \\ \midrule
\TODO{họ tên} & \TODO{} & \TODO{} \\ \TODO{họ tên} & \TODO{} & \TODO{} \\ \bottomrule\end{tabular}

\section{Mục lục, nguồn dữ liệu và tài liệu tham khảo}
\subsection*{Nguồn dữ liệu}
Sáu bộ dữ liệu thuộc gói R \texttt{datasets}, tải từ Rdatasets (\url{https://vincentarelbundock.github.io/Rdatasets/}), ngày tải 20/11/2026. \TODO{chép trích dẫn nguyên văn từ \texttt{?LakeHuron}, \texttt{?lynx}, \texttt{?EuStockMarkets}, \texttt{?Nile}, \texttt{?AirPassengers}.}
\subsection*{Tài liệu tham khảo}
\begin{enumerate}
\item Dung Nguyen, AS3179 Mô hình thống kê, Bài 11: Chuỗi thời gian I -- Phụ thuộc và mô hình ARMA; Bài 12: Chuỗi thời gian II -- ARIMA và dự báo (slide bài giảng). Khoa Khoa học Ứng dụng, HCMUT. \TODO{ghi năm học}
\item Box, G.E.P., Jenkins, G.M., Reinsel, G.C., Ljung, G.M., \textit{Time Series Analysis: Forecasting and Control}, 5th ed., Wiley, 2015.
\item Brockwell, P.J., Davis, R.A., \textit{Introduction to Time Series and Forecasting}, 3rd ed., Springer, 2016.
\item Shumway, R.H., Stoffer, D.S., \textit{Time Series Analysis and Its Applications}, 4th ed., Springer, 2017.
\item Hyndman, R.J., Athanasopoulos, G., \textit{Forecasting: Principles and Practice}, 3rd ed., OTexts.
\item Box, G.E.P., Pierce, D.A. (1970), JASA 65. \quad Ljung, G.M., Box, G.E.P. (1978), Biometrika 65. \quad Bartlett, M.S. (1946), J. Roy. Statist. Soc. Suppl. 8.
\item Dickey, D.A., Fuller, W.A. (1979), JASA 74. \quad Kwiatkowski, Phillips, Schmidt, Shin (1992), J. Econometrics 54. \quad Diebold, F.X., Mariano, R.S. (1995), J. Bus. Econ. Statist. 13.
\item Hyndman, R.J., Koehler, A.B. (2006), Int. J. Forecasting 22 (MASE). \quad Engle, R.F., Granger, C.W.J. (1987), Econometrica 55.
\item Seabold, S., Perktold, J. (2010), Statsmodels: Econometric and statistical modeling with Python, Proc. 9th Python in Science Conf.
\end{enumerate}
\TODO{đối chiếu lại số trang/tập của từng tài liệu với danh sách ở cuối hai bài giảng; chỉ giữ những tài liệu thực sự được dùng.}

\end{document}
